\documentclass[conference]{IEEEtran}

\usepackage{cite}
\usepackage{amsmath,amssymb,amsfonts}
\usepackage{graphicx}
\usepackage{textcomp}
\usepackage{xcolor}
\usepackage{booktabs}
\usepackage{tabularx}
\usepackage{array}
\usepackage{url}

% Arrow definitions may be placed in the preamble.
\newcommand{\dn}{\,\ensuremath{\downarrow}}
\newcommand{\up}{\,\ensuremath{\uparrow}}
\newcommand{\same}{\,\ensuremath{\leftrightarrow}}
\newcommand{\dnsig}{\,\ensuremath{\downarrow^{*}}}
\newcommand{\upsig}{\,\ensuremath{\uparrow^{*}}}
\newcommand{\dnns}{\,\ensuremath{\downarrow^{\mathrm{ns}}}}
\newcommand{\upns}{\,\ensuremath{\uparrow^{\mathrm{ns}}}}

\newcolumntype{Y}{>{\centering\arraybackslash}X}

\raggedbottom

\begin{document}

\title{ASR-Style Transcript Normalization Removes Useful Pragmatic Signal for Dialogue Act Classification}

\author{
\IEEEauthorblockN{Lucky Onyekwelu-Udoka}
\IEEEauthorblockA{\textit{Department of Electrical and Computer Engineering} \\
\textit{Iowa State University}\\
Ames, Iowa, United States \\
lucky@iastate.edu}
\and
\IEEEauthorblockN{Md Shafiqul Islam}
\IEEEauthorblockA{\textit{Department of Electrical and Computer Engineering} \\
\textit{Iowa State University}\\
Ames, Iowa, United States \\
shafiqul@iastate.edu}
}

\maketitle

\begin{abstract}
Automatic speech recognition (ASR) pipelines often
normalize transcripts by removing filled pauses,
punctuation, and colloquial expressions, potentially
discarding cues useful for dialogue act classification.
We evaluate nine transcript conditions across six
classifiers (TF-IDF logistic regression, RoBERTa,
BERT, ALBERT, GPT-2, and BiLSTM) and three conversational
corpora (Switchboard, ICSI MRDA, and SLUE-HVB),
supplemented by multi-seed experiments, normalization
controls, and real-ASR evaluations. Across five random
seeds, heavy normalization reduces Macro-F1 in all
15 neural model--corpus combinations under the original
preprocessing protocols. Removal with the original filler/backchannel lexicon
consistently degrades performance on MRDA and SLUE-HVB
but produces mixed effects on Switchboard. Punctuation-aware TF-IDF
experiments show that punctuation removal reduces
Macro-F1 by 0.036 on Switchboard and 0.078 on MRDA,
revealing that the original linear baseline's
insensitivity was a tokenization artifact. Restricting
filler removal to \textit{um}, \textit{uh}, and
\textit{erm} preserves degradation on MRDA but largely
eliminates the SLUE-HVB effect, where the original
lexicon removes class-defining backchannel expressions.
These findings demonstrate that normalization effects
depend on preprocessing definitions, feature extraction,
and conversational context, supporting task-dependent
rather than indiscriminate transcript normalization.
\end{abstract}

\begin{IEEEkeywords}
automatic speech recognition, transcript normalization,
dialogue act classification, disfluency, filled pauses
\end{IEEEkeywords}

\section{Introduction}

Automatic speech recognition (ASR) transcripts are
often normalized by removing filled pauses, stripping
punctuation, and standardizing colloquial expressions.
Although these operations improve readability, they
may discard cues relevant to turn management and
discourse structure
\cite{CLARK200273,
https://doi.org/10.1111/j.1749-818X.2008.00068.x,
CLARK1989259}. While previous studies have examined
disfluencies and ASR errors, how normalization affects
dialogue act classification across models and
conversational settings remains insufficiently
characterized.

We investigate nine transcript conditions across
six classifiers (TF-IDF logistic regression, RoBERTa,
BERT, ALBERT, GPT-2, and BiLSTM) and three corpora
(Switchboard, ICSI MRDA, and SLUE-HVB), yielding
18 model--corpus combinations. The main experiments
transform reference transcripts; separate analyses
examine recognition errors in real ASR output.
We also assess tokenization, lexicon sensitivity,
and random-seed variability.

Our main contributions are:
\begin{enumerate}
    \item A multi-model, cross-corpus evaluation
    showing consistent Macro-F1 degradation under
    heavy normalization across five random seeds,
    with more variable effects from filler removal.

    \item A punctuation-aware TF-IDF analysis
    demonstrating that the apparent insensitivity
    of linear classifiers to punctuation removal
    results from tokenization rather than model
    architecture alone.

    \item A restricted filler-lexicon analysis
    showing that degradation persists on MRDA
    but largely disappears on SLUE-HVB when
    class-defining backchannel tokens are preserved.

    \item Cluster-level statistical analyses,
    real-ASR evaluations, and selective
    normalization experiments assessing the
    robustness and practical implications
    of transcript preprocessing.
\end{enumerate}

Our findings show that normalization effects depend
on the linguistic cues removed, feature extraction,
and conversational setting, supporting task-dependent
rather than indiscriminate transcript preprocessing.

\section{Related Work}

Disfluencies, filled pauses, and backchannels
convey information about hesitation, turn management,
and discourse structure
\cite{CLARK200273,
https://doi.org/10.1111/j.1749-818X.2008.00068.x,
shriberg1999phoneticCO,
GinzburgFernandezSchlangen2014,
CLARK1989259}. Removing these expressions during
transcript normalization may therefore eliminate
cues useful for downstream pragmatic inference.

Prior studies have investigated ASR recognition errors
and configurable transcription-noise evaluation
\cite{Cui2021AnAT,Nechaev2021,LI2024104598,
kubis-etal-2023-back,shapira-etal-2025-measuring,
zusag24_interspeech}.
However, recognition errors differ from deliberate
transcript normalization, which can alter linguistic
cues even when transcription is accurate.

Dialogue act recognition research has examined
lexical and conversational cues
\cite{jurafsky-shriberg-biasca1997,
Core1997CodingDW,stolcke-etal-2000-dialogue}.
Zelasko et al.\ \cite{zelasko-etal-2021-helps}
investigated the role of punctuation in
transformer-based dialogue act recognition and
segmentation.

Disfluency studies establish the communicative role
of hesitations and backchannels, whereas ASR-robustness
work examines recognition errors and Zelasko et al.\
focus on punctuation in dialogue-act recognition.
These studies do not jointly isolate deliberate
filler/backchannel removal, punctuation stripping,
and colloquial expansion for dialogue-act classification
across model families and conversational registers.
Our controlled text ablations address this gap with
nine conditions, six classifiers, and three corpora,
supplemented by tokenizer and lexicon controls,
multi-seed evaluation, and real-ASR experiments.



\section{Experimental Setup}

\textbf{Datasets.}
We evaluate three conversational corpora: Switchboard
\cite{chang-et-al2020}, ICSI MRDA
\cite{shriberg-et-al2004-mrda,wylupek-mrda},
and SLUE-HVB \cite{shon-et-al2023-slue2},
representing telephone conversations, multi-party
meetings, and scripted call-center dialogues,
respectively. SLUE-HVB is multi-label; approximately
40\% of utterances have multiple tags, and we use the
release's first tag as the primary class. Table~\ref{tab:datasets} summarizes
the datasets after removing labels with fewer than
50 instances in combined train--test counts. Switchboard uses a conversation-level
80/20 split (seed 42), while MRDA and HVB use their
release-defined splits.

\begin{table}[t]
\caption{Dataset statistics after rare-label filtering
($\geq 50$ occurrences).}
\label{tab:datasets}
\centering
\scriptsize
\setlength{\tabcolsep}{3pt}
\renewcommand{\arraystretch}{0.9}
\begin{tabularx}{0.95\columnwidth}{@{}lYYY@{}}
\toprule
& \textbf{Switchboard} & \textbf{MRDA} & \textbf{HVB} \\
\midrule
Register & Telephone & Meetings & Call-center \\
Train/test & 34.6k/9.1k & 75.1k/16.7k & 11.3k/6.1k \\
Labels & 27 & 12 & 16 \\
Canonical filler rate$^{a}$ & 40.5\% & 15.9\% & 5.7\% \\
\bottomrule
\end{tabularx}
\vspace{1pt}\parbox{0.95\columnwidth}{\tiny $^{a}$Utterances containing \textit{um}, \textit{uh}, or \textit{erm}.}
\end{table}

\textbf{Models.}
We evaluate TF-IDF with logistic regression
(TF-IDF+LR), RoBERTa-base \cite{liu-et-al2019-roberta},
BERT-base-cased \cite{devlin-et-al2018-bert},
ALBERT-base-v2 \cite{lan-et-al2019-albert},
GPT-2 \cite{radford-et-al2019-gpt2}, and BiLSTM
\cite{hochreiter-schmidhuber1997-lstm}.
TF-IDF+LR uses unigram/bigram features
($\mathrm{min\_df}=3$, $\mathrm{max\_df}=0.95$)
and class-balanced logistic regression.
Transformers are fine-tuned for three epochs
(learning rate $2\times10^{-5}$, batch size 16,
maximum length 96). BiLSTM uses 200-dimensional
embeddings, 128 hidden units, and 15 training
epochs. Neural experiments use five seeds
(1, 2, 3, 4, and 42); TF-IDF+LR is deterministic.
We additionally evaluate a punctuation-aware,
case-preserving TF-IDF tokenizer to examine
punctuation discarded by the default tokenizer.

\textbf{Normalization Conditions.}
We evaluate nine conditions:
\texttt{gold\_text} (original transcripts),
\texttt{norm\_colloq\_only} (colloquial expansion),
\texttt{rm\_punct\_only} (punctuation removal),
\texttt{rm\_discourse\_only} (removal of
\textit{like}/\textit{well}),
\texttt{gold\_alpha\_text} (corpus-specific alphanumeric form),
\texttt{rm\_fillers\_only} (filler/backchannel lexicon removal),
\texttt{clean\_light} (lowercasing and removal
of non-alphanumeric characters except apostrophes),
\texttt{clean\_heavy} (light cleaning, filler
removal, and colloquial expansion), and
\texttt{asr\_noise} (synthetic recognition errors).

Synthetic ASR noise applies token deletion (8\%),
substitution (8\%), and insertion (4\%).
Each main-grid condition is applied to both
training and test transcripts.

The original filler lexicon contains
\textit{um}, \textit{uh}, \textit{erm},
\textit{mm}, \textit{hmm}, \textit{mhm},
\textit{uhhuh}, \textit{uh-huh}, and
\textit{huh}. The restricted lexicon retains
only \textit{um}, \textit{uh}, and \textit{erm}.
A separate punctuation--filler factorial control
evaluates removal of each feature and both together
using punctuation-aware TF-IDF+LR on MRDA.

An additional \texttt{pragmatic\_norm} control
applies lowercasing and colloquial expansion
while preserving fillers, backchannels, and
punctuation. This transformation leaves MRDA
unchanged; its reported neural comparisons
are single-seed supplementary controls.

Normalization implementations have corpus-specific
differences. Switchboard's original
\texttt{clean\_heavy} additionally removes
\textit{like} and \textit{youknow}, while its \texttt{gold\_alpha\_text}
uses ConvoKit's provided field, which also removes
clitics. MRDA and HVB instead use regex-based
alphanumeric filtering. Colloquial normalization
changes no Switchboard or MRDA utterances, while
punctuation removal and light cleaning leave HVB
unchanged.

As a separate single-seed sensitivity test, we
evaluate a harmonized strict heavy-normalization
condition applying
lowercasing, punctuation removal, colloquial
expansion, and deletion of only \textit{um},
\textit{uh}, and \textit{erm} across all
three corpora.

\textbf{Real-ASR Evaluation.}
We evaluate Whisper-base.en
\cite{radford-et-al2022-whisper} under two protocols.
On SLUE-HVB, models are trained and tested on
ASR transcripts (17.7\%/18.3\% train/test WER)
and compared with gold-trained/gold-tested models.
On MRDA, gold-trained models are compared on
gold and ASR test transcripts from a separately
reconstructed 60/15 meeting-disjoint split
(41.0\% WER), combining recognition errors
with train--test transcription mismatch.
Switchboard uses synthetic ASR noise because
its audio requires an LDC license.

\textbf{Evaluation.}
Macro-F1 is the primary evaluation metric,
assigning equal weight to each dialogue-act class.
Primary neural results are summarized using
mean and standard deviation across five seeds.
At the primary seed (42), paired Macro-F1
differences are assessed using 2,000
cluster-bootstrap resamples, with conversations
as clusters for Switchboard and calls for HVB.
McNemar's test assesses paired utterance-level
correctness rather than Macro-F1, with
Bonferroni correction across 36 comparisons; 21
remain significant at $\alpha=0.05/36$.
Because utterances within conversations are
dependent, McNemar results are supplementary
to cluster-level uncertainty estimates.

Implementation scripts and supplementary
materials are provided in a public project
repository.\footnote{
\url{https://github.com/luckymaduabuchi/asr-normalization-dialogue-acts}}




\section{Results}
\label{sec:results}

\begin{table*}[t]
\caption{Macro-F1 under original, filler-removed, and
heavily normalized transcripts. Neural results are
mean $\pm$ standard deviation over five seeds.
TF-IDF+LR is deterministic. Original corpus-specific
normalization protocols are used.}
\label{tab:grid}
\centering
\footnotesize
\setlength{\tabcolsep}{5pt}
\renewcommand{\arraystretch}{0.95}
\begin{tabular*}{\textwidth}
{@{\extracolsep{\fill}}llccc@{}}
\toprule
\textbf{Corpus} & \textbf{Model} &
\textbf{Gold F1} & \textbf{Filler F1} &
\textbf{Heavy F1} \\
\midrule
Switchboard & TF-IDF+LR & 0.311 & 0.300 & 0.293 \\
            & RoBERTa & $0.449\pm0.006$ & $0.447\pm0.005$ & $0.393\pm0.012$ \\
            & BERT    & $0.440\pm0.005$ & $0.449\pm0.008$ & $0.387\pm0.004$ \\
            & ALBERT  & $0.420\pm0.008$ & $0.409\pm0.026$ & $0.355\pm0.046$ \\
            & GPT-2   & $0.388\pm0.012$ & $0.383\pm0.009$ & $0.339\pm0.008$ \\
            & BiLSTM  & $0.370\pm0.007$ & $0.374\pm0.012$ & $0.336\pm0.010$ \\
\midrule
MRDA        & TF-IDF+LR & 0.409 & 0.364 & 0.361 \\
            & RoBERTa & $0.593\pm0.006$ & $0.550\pm0.003$ & $0.485\pm0.003$ \\
            & BERT    & $0.591\pm0.006$ & $0.556\pm0.002$ & $0.481\pm0.007$ \\
            & ALBERT  & $0.580\pm0.007$ & $0.530\pm0.016$ & $0.469\pm0.013$ \\
            & GPT-2   & $0.565\pm0.001$ & $0.522\pm0.008$ & $0.444\pm0.003$ \\
            & BiLSTM  & $0.496\pm0.005$ & $0.438\pm0.006$ & $0.401\pm0.011$ \\
\midrule
SLUE-HVB    & TF-IDF+LR & 0.713 & 0.660 & 0.661 \\
            & RoBERTa & $0.745\pm0.008$ & $0.697\pm0.004$ & $0.698\pm0.004$ \\
            & BERT    & $0.742\pm0.004$ & $0.698\pm0.004$ & $0.697\pm0.003$ \\
            & ALBERT  & $0.742\pm0.004$ & $0.697\pm0.003$ & $0.698\pm0.003$ \\
            & GPT-2   & $0.707\pm0.026$ & $0.671\pm0.004$ & $0.673\pm0.007$ \\
            & BiLSTM  & $0.732\pm0.005$ & $0.684\pm0.008$ & $0.683\pm0.011$ \\
\bottomrule
\end{tabular*}
\end{table*}


\subsection{Multi-Seed Normalization Effects}

Table~\ref{tab:grid} reports Macro-F1 across five seeds
for the neural models and a deterministic TF-IDF+LR
baseline. Under the original preprocessing protocols,
heavy normalization reduces Macro-F1 in all 75 paired
seed--model--corpus comparisons. Filler removal also
reduces Macro-F1 in all 50 neural comparisons on MRDA
and SLUE-HVB but produces mixed effects on Switchboard.
For example, MRDA RoBERTa decreases from 0.593 to
0.550 after filler removal and to 0.485 under heavy
normalization. In contrast, Switchboard BERT improves
from 0.440 to 0.449 after filler removal but falls
to 0.387 under heavy normalization. Thus, heavy
normalization produces more consistent degradation
than filler removal alone.


\subsection{Punctuation and Filler-Lexicon Controls}

The original TF-IDF+LR baseline appeared insensitive
to punctuation removal because its default tokenizer
discarded punctuation before feature extraction.
Using a punctuation-aware, case-preserving tokenizer,
punctuation removal reduces Macro-F1 by 0.036 on
Switchboard and 0.078 on MRDA (Table~\ref{tab:punct}).
Thus, the original near-neutral effect was a
tokenization artifact rather than an intrinsic
property of linear classifiers. SLUE-HVB remains
unaffected because its transcripts lack punctuation.

\begin{table}[t]
\caption{Punctuation-removal effects for TF-IDF+LR.
$\Delta_{\mathrm{punct}}$ denotes the Macro-F1 decrease
relative to the corresponding gold condition.}
\label{tab:punct}
\centering
\scriptsize
\setlength{\tabcolsep}{4pt}
\begin{tabular}{@{}lccc@{}}
\toprule
\textbf{Corpus} &
\textbf{Default $\Delta$} &
\textbf{Aware Gold} &
\textbf{Aware $\Delta$} \\
\midrule
Switchboard & 0.0002 & 0.360 & 0.036 \\
MRDA        & 0.0000 & 0.494 & 0.078 \\
SLUE-HVB    & 0.0000 & 0.717 & 0.000 \\
\bottomrule
\end{tabular}
\end{table}

We also evaluate a restricted filler lexicon containing
only \textit{um}, \textit{uh}, and \textit{erm}.
On MRDA, filler-removal degradation remains positive
across all six models ($\Delta=0.015$--$0.039$),
although the proportion of emptied backchannel
utterances falls from 43.1\% to 0.5\%.
Across the filtered SLUE-HVB data, the original
lexicon empties 38 of 50 backchannel utterances,
whereas the restricted lexicon empties none; only
16 backchannels occur in the held-out test set. Accordingly, filler-removal degradation
approaches zero for the four evaluated neural models
with non-anomalous primary-seed baselines
($\Delta=-0.006$ to $0.001$).

\begin{table}[t]
\caption{Primary-seed Macro-F1 degradation under
harmonized strict heavy normalization.
$\Delta=\mathrm{gold}-\mathrm{strict}$;
positive values indicate degradation.}
\label{tab:strict}
\centering
\scriptsize
\setlength{\tabcolsep}{5pt}
\renewcommand{\arraystretch}{1.0}
\begin{tabular}{@{}lrrr@{}}
\toprule
\textbf{Model} &
\textbf{SWB} & \textbf{MRDA} & \textbf{HVB} \\
\midrule
TF-IDF+LR & 0.007 & 0.039 & 0.015 \\
RoBERTa   & 0.044 & 0.080 & $-0.004$ \\
BERT      & 0.041 & 0.076 & $-0.001$ \\
ALBERT    & 0.032 & 0.077 & 0.003 \\
GPT-2     & 0.048 & 0.110 & $-0.049$ \\
BiLSTM    & 0.022 & 0.056 & 0.000 \\
\bottomrule
\end{tabular}
\end{table}

Table~\ref{tab:strict} reports results under an
identical strict heavy-normalization function.
All six models exhibit Macro-F1 degradation on
Switchboard and MRDA, whereas five of six SLUE-HVB models change
by at most 0.015. The exception is GPT-2, which
improves by 0.049 from an unusually low gold
baseline. All strict-control results use seed 42.

Separately, omitting Switchboard's extra
\textit{like}/\textit{youknow} removal while retaining the
nine-word lexicon changes the normalized text
for 17.7\% of utterances. TF-IDF+LR Macro-F1
changes from 0.2933 to 0.2954; corresponding
neural reruns are unavailable.

At seed 42, \texttt{pragmatic\_norm} Macro-F1
remains within $\pm0.013$ of gold across 12
Switchboard and SLUE-HVB model combinations.
It outperforms heavy normalization in 11 of 12
comparisons, suggesting that mild standardization
may preserve performance when discourse cues remain.

\subsection{Statistical and Class-Level Analysis}

Five-seed standard deviations summarize training
variability. Paired cluster-bootstrap confidence
intervals quantify test-cluster sampling uncertainty
at seed 42, not uncertainty in five-seed means.
Throughout, positive $\Delta$ denotes the
Macro-F1 decrease relative to gold.

At the primary seed, cluster-bootstrap 95\% confidence
intervals for heavy-normalization Macro-F1 differences
exclude zero for all five neural models on Switchboard
and four of five on SLUE-HVB. Filler-removal intervals
exclude zero for only ALBERT on Switchboard and four
neural models on SLUE-HVB. For TF-IDF+LR on Switchboard,
filler-removal and heavy-normalization differences are
0.011 $[0.003,0.020]$ and 0.017 $[0.009,0.027]$,
respectively. These are nominal 95\% intervals without
multiplicity adjustment; an F1 decrease alone does
not establish statistical significance. The primary
MRDA grid lacks recoverable meeting identifiers. On a
separate 60/15 meeting-disjoint reconstruction, heavy-
normalization intervals exclude zero for four of five
neural models, while filler-only effects are split-sensitive.

Utterance-level McNemar tests provide an additional
paired-correctness comparison but are not tests
of Macro-F1 and do not adjust for conversation
clustering; cluster-bootstrap intervals therefore
remain the primary uncertainty evidence.

\begin{table}[t]
\caption{Selected class-level effects of lexicon-based
removal for deterministic TF-IDF+LR on held-out test
splits. P, R, and F1 denote precision, recall, and
F1-score. BC denotes SLUE-HVB backchannels.}
\label{tab:class-effects}
\centering
\scriptsize
\setlength{\tabcolsep}{4pt}
\renewcommand{\arraystretch}{1.05}
\begin{tabular}{@{}llcrr@{}}
\toprule
\textbf{Corpus} & \textbf{Class} &
\textbf{Metric} & \textbf{Gold} & \textbf{No Fill.} \\
\midrule
Switchboard & b  & P  & 0.558 & 0.436 \\
            & b  & F1 & 0.433 & 0.365 \\
\midrule
MRDA        & fh & R  & 0.495 & 0.326 \\
            & h  & R  & 0.313 & 0.104 \\
\midrule
SLUE-HVB    & BC & P  & 0.353 & 0.092 \\
            & BC & R  & 0.750 & 0.750 \\
            & BC & F1 & 0.480 & 0.163 \\
\bottomrule
\end{tabular}
\end{table}

Table~\ref{tab:class-effects} shows different
class-level effects across corpora. On Switchboard,
backchannel precision decreases from 0.558 to 0.436,
with an F1 reduction of 0.068. On MRDA, recall
decreases for floor holders (\texttt{fh}) and holds
before answers or agreements (\texttt{h}),
indicating difficulty recognizing turn-management
dialogue acts. In contrast, SLUE-HVB backchannel
recall remains at 0.750, while precision decreases
from 0.353 to 0.092, reducing F1 from 0.480 to
0.163. Because Macro-F1 weights this low-support class
equally with the other 15 classes, its large F1
change warrants cautious interpretation.
The restricted-lexicon analysis indicates that
much of this degradation results from deleting
class-defining backchannel expressions.

\subsection{Real-ASR Evaluation}

Whisper-base.en is evaluated under two protocols.
On SLUE-HVB, models are trained and tested on
Whisper transcripts (17.7\%/18.3\% train/test WER).
Five models exhibit lower Macro-F1, whereas
GPT-2 improves from 0.6674 to 0.6765 at seed 42.
This isolated result does not establish an ASR
benefit, especially given its unusually low baseline.

On MRDA, models trained on gold transcripts are
evaluated on Whisper transcripts from 15 held-out
meetings (WER: 41.0\%). All six models degrade,
with Macro-F1 decreases ranging from 0.1366 to
0.2822. Unlike SLUE-HVB, this experiment combines
recognition errors with train--test transcription
mismatch; the effect sizes are therefore not directly
comparable.


\section{Discussion and Conclusion}
\label{sec:conclusion}

Our findings show that transcript normalization affects
dialogue act classification differently across models,
corpora, and preprocessing operations. Under the original
protocols, heavy normalization reduces Macro-F1 across
all five seeds for every neural model--corpus combination,
whereas filler removal produces mixed effects on
Switchboard but consistent degradation on MRDA and
SLUE-HVB.

\textbf{Interpretation of normalization effects.}
The punctuation-aware TF-IDF control shows
that punctuation is predictive when the feature
extractor preserves it. The default TF-IDF
baseline appeared insensitive because its tokenizer
discarded punctuation. Filler-lexicon controls
reveal a different mechanism: on MRDA, degradation
persists when removal is restricted to \textit{um},
\textit{uh}, and \textit{erm}, suggesting that these
expressions provide useful classification cues.
On SLUE-HVB, most degradation disappears under the
restricted lexicon because the original inventory
removes class-defining backchannel expressions.
Thus, not all performance losses reflect the same
type of information removal.

\textbf{Robustness and practical implications.}
\label{sec:swb-recovery}
The five-seed experiments show that heavy-normalization
degradation is more consistent than filler-removal
effects. Under harmonized strict normalization,
degradation persists across all six classifiers on
Switchboard and MRDA but is largely absent on
SLUE-HVB, demonstrating sensitivity to corpus
characteristics and lexicon definitions.
Single-seed pragmatic-normalization controls
further suggest that modest text standardization
can retain dialogue-act performance when
backchannels and punctuation are preserved.

\textbf{Limitations.}
The study is limited to three English-language
corpora. The original Switchboard heavy-normalization
implementation additionally removes \textit{like} and \textit{youknow},
and harmonized lexicon experiments use a single seed.
On Switchboard, spaced backchannel expressions such
as \textit{Uh - huh} are partially affected by the
strict lexicon, which removes \textit{uh} but retains
\textit{huh}. Thus, identical normalization rules
do not necessarily remove equivalent linguistic cues
across corpora. The primary MRDA partition is
meeting-disjoint but not speaker-disjoint, and
meeting-level confidence intervals are unavailable
for that primary split. The alternative MRDA
neural comparisons use different training scripts,
hyperparameters, and text formatting, preventing
controlled attribution of their differences to
normalization. Real-ASR evaluations use different
protocols: matched ASR training/testing on SLUE-HVB
and gold-trained/ASR-tested models on MRDA.
The pragmatics-preserving control uses reference
transcripts and does not recover cues already
missing from real ASR output. All reported BiLSTM
results use the corrected empty-input pooling
implementation.

Overall, aggressive transcript normalization can
degrade dialogue act classification, but the
magnitude and mechanism depend on feature extraction,
lexicon design, and conversational setting.
These findings support selective, task-dependent
normalization rather than indiscriminate transcript
cleaning. Future work will examine multi-seed harmonized
normalization, matched ASR training on MRDA, and
additional conversational domains.

\bibliographystyle{IEEEtran}
\bibliography{asilomar_refs}

\end{document}
